\section{Recursão nos números naturais}

Intuitivamente, o significado de $m+n$ consiste em tomar $m$ e aplicar a operação de sucessor $n$ vezes.
Como uma criança aprendendo a somar, parte-se de $m$ e, ``contando nos dedos'', percorre-se $n$ números, e o resultante é $m+n$.
Assim, $m+5$, por exemplo, deveria ser definido como $S(S(S(S(S(m)))))$.
Pensamos em $m+0$ como o próprio $m$.
De modo geral, pensa-se que $m+n=S(S(\dots S(m)\dots))$, onde $S$ é aplicado $n$ vezes.
Porém, como formalizar isso?

Uma definição recursiva que costuma ser apresentada é a seguinte.

Fixado $m \in \omega$, quer-se definir $m+n$ \emph{recursivamente}, de modo que:

\begin{equation}\label{eq:recursao-soma}
    m+n =
    \begin{cases}
        m & \text{se } n=0\\
        S(m+k) & \text{se } n=S(k)
    \end{cases}
\end{equation}
Seguindo o roteiro dado pela Equação~\ref{eq:recursao-soma}, temos:

\begin{align*}
m+0 &= m\\
m+1 &= m+S(0) = S(m+0) = S(m)\\
m+2 &= m+S(1) = S(m+1) = S(S(m))\\
m+3 &= m+S(2) = S(m+2) = S(S(S(m)))\\
m+4 &= m+S(3) = S(m+3) = S(S(S(S(m))))\\
\end{align*}
\emph{E assim por diante.}

Parece funcionar. Mas podemos, de fato, definir $m+n$ por meio da Equação~\ref{eq:recursao-soma}?
Tal equação é, de certa forma, autorreferente: usa-se $+$, ainda que para ``números menores'', para definir $+$.
Lembramos ao leitor que construções autorreferentes, no geral, podem ser problemáticas em Teoria dos Conjuntos (ver, por exemplo, o Exercício~\ref{exer:compreensao-autorreferente}).
Assim, faz-se necessária uma justificativa formal apropriada para a definição dada pela Equação~\ref{eq:recursao-soma} (que, conforme veremos, será, sim, lícita).
Para chegarmos ao enunciado do Teorema da recursão, vamos investigar o que está acontecendo na Equação~\ref{eq:recursao-soma}.

Lá, temos um \emph{parâmetro} fixado, $m$.
Podemos pensar que estamos, para cada $m$, definindo uma sequência $f_m$ de números naturais, em que $f_m(n)$ representa $m+n$.

Além disso, para definirmos $f_m(S(k))$, precisamos conhecer $f_m(k)$.
Assim, a definição de $f_m(n)$ é dada por uma condição que envolve um parâmetro $m$ e o valor de $f_m$ para o antecessor de $n$ (caso $n\neq 0$).

Porém, apenas o conhecimento do valor anterior não é sempre suficiente.
Temos, por exemplo, a seguinte famosa sequência de números naturais, conhecida como \emph{sequência de Fibonacci}:
\begin{align*}
F(0) &= 0\\
F(1)&= 1\\
F(n)&= F(n-1) + F(n-2) \text{ para } n\geq 2
\end{align*}

Fica claro que, para $n\geq 2$ $F(n)$ não depende apenas do valor anterior, mas dos dois anteriores.

A exposição recursiva da sequência de Fibonacci acima parte de $0$ e $1$ como primeiros elementos.
Algumas pessoas preferem iniciá-la de $1$ e $1$.
No geral, poderíamos utilizar qualquer par de números para iniciar a sequência (e obter assim sequências potencialmente diferentes).
Podemos pensar no par $(0, 1)$ como um parâmetro da sequência.
No geral, podemos pensar em definir uma sequência de Fibonacci parametrizada por um par $p=(a, b)$ de números naturais, em que $F_p(0)=a$, $F_p(1)=b$ e $F_p(n)=F_p(n-1)+F_p(n-2)$ para $n\geq 2$.



No geral, ao se definir uma sequência $s_p=(s_p(n): n \in \omega)$, $s_n$ pode depender de todos os valores anteriores, $t=(s_p(k): k < n)$, e não apenas do valor imediatamente anterior, $s_{n-1}$, bem como de parâmetros $p$ que ajam em toda a sequência, como $m$ na definição de $m+n$.
Definida uma sequência inicial $t$ da sequência final, efetuamos uma operação $G(p, t)$ para obter o próximo elemento da sequência.

No caso, $G:\omega^2\times \omega^{<\omega}\to \omega$ é dada por:

\begin{align*}
G(p, t) =
\begin{cases}
    p(0) & \text{se } t=\emptyset\\
    p(1) & \text{se } t \in \omega^1\\
    t(n-1)+t(n-2) & \text{se } t\in \omega^n \text{ para } n\geq 2
\end{cases}
\end{align*}

Na prática, como calcular $F_p(n)$?
Ora, a priori, deve-se calcular toda a sequência $F_p(0), F_p(1), \dots, F_p(n-1)$, de forma iterada, até chegar em $F_p(n)$.
Iniciamos sem conhecimento de nenhum elemento, ou seja, temos a sequência vazia $t_0=\emptyset$.
Calcula-se $F_p(0)=G(p, t_0)=p(0)$, e, assim, $t_1=(F_p(0))$.
Calcula-se $F_p(1)=G(p, t_1)=p(1)$, e, assim, $t_2=(F_p(0), F_p(1))$.
Calcula-se $F_p(2)=G(p, t_2)=t_2(1)+t_2(0)=F_p(1)+F_p(0)$, e, assim, $t_3=(F_p(0), F_p(1), F_p(2))$.
Procedendo-se dessa forma, chega-se a $F_p(n)=G(p, t_n)$, onde $t_n=(F_p(k): k<n)$.
Tal raciocínio é formalizado na prova do Teorema da Recursão, enunciado abaixo.

Antes, introduziremos uma notação.

\begin{definition}
    Sejam $A, B, C$ conjuntos.
    Dada $f: A\times B\to C$ e $a\in A$, definimos $f_a: B\to C$ por $f_a(b)=f(a, b)$.

    Reciprocamente, dada uma família de funções $f_a: B\to C$ indexada por $a\in A$, definimos $f: A\times B\to C$ por $f(a, b)=f_a(b)$.
\end{definition}

\begin{theorem}[Teorema da Recursão]\label{thm:recursion}
    Sejam $A, P$ conjuntos e $G: P\times A^{<\omega}\to A$ uma função.
    Então existe uma única função $f:P\times \omega\to A$ tal que, para todo $p \in P$ e $n\in \omega$, $f_p(n)=G(p, f_p\restriction n)$.
\end{theorem}

\begin{proof}
    Dado $p \in P$ e $n \in \omega$, define-se que $t \in A^{<\omega}$ é uma \emph{computação parcial baseada em $G$ e $p$ de comprimento $n$} se, e somente se, $t\in A^n$ e, para todo $k<n$, $t(k)=G(p, t\restriction k)$.
    
    \begin{claim}
        Se $p \in P$ e $t$ e $t'$ são computações parciais baseadas em $G$ e $p$ de comprimento $n$, $n'$ respectivamente, então $t$ e $t'$ são compatíveis.
    \end{claim}
    \begin{proof}[Prova da afirmação]
        Sem perda de generalidade, suponha que $n\leq n'$.
        Assim, $\dom t\cap \dom t'=\dom t=n$.
        Devemos ver que para todo $k<n$, $t(k)=t'(k)$.

        Suponha que não.
        Seja $m$ o menor número natural tal que $m<n$ e $t(m)\neq t'(m)$.
        Assim, $t\restriction m=t'\restriction m$.
        Logo, $t(m)=G(p, t\restriction m)=G(p, t'\restriction m)=t'(m)$, o que é uma contradição.
    \end{proof}
    \begin{claim}
        Para todo $p \in  P$ e $n \in \omega$, existe uma computação parcial baseada em $G$ e $p$ de comprimento $n$.
    \end{claim}
    \begin{proof}[Prova da afirmação]
        Procede-se por indução em $n$.
        Para $n=0$, a sequência vazia é uma computação parcial baseada em $G$ e $p$ de comprimento $0$.

        Suponha que $n\in \omega$ e que existe uma computação parcial baseada em $G$ e $p$ de comprimento $n$.
        Seja $t$ tal computação.
        Assim, $t\in A^n$ e, para todo $k<n$, $t(k)=G(p, t\restriction k)$.
        Defina-se $t'=t\cup \{(n, G(p, t))\}$.
        Assim, $t'\in A^{n+1}$ e, para todo $k<n+1$, temos que:
        \begin{itemize}
            \item se $k<n$, então $t'(k)=t(k)=G(p, t\restriction k)=G(p, t'\restriction k)$;
            \item se $k=n$, então $t'(k)=G(p, t)=G(p, t'\restriction k)$.
        \end{itemize}
        Logo, $t'$ é uma computação parcial baseada em $G$ e $p$ de comprimento $n+1$, o que completa a indução.
    \end{proof}

    Dado $p \in P$, seja $f_p=\bigcup\{t: t \text{ é uma computação parcial baseada em } G \text{ e } p\}$.

    Formalmente, acima estamos definindo $h:P\to\mathcal P(\omega\times A)$ por
    \[h(p)=\bigcup \{t \in [\omega\times A]^{<\omega}: t \text{ é uma computação parcial baseada em } G \text{ e } p\},
    \]
    o que é possível pelo Esquema da Compreensão, e escrevendo $f_p=h(p)$.

    Para cada $p \in P$, $f_p$ é uma função, pois é uma união de funções compatíveis.
    Além disso, $\dom f_p = \bigcup\{\dom t: t \text{ é uma computação parcial baseada em } G \text{ e } p\}=\omega$, pois, para cada $n\in \omega$, existe uma computação parcial baseada em $G$ e $p$ de comprimento $n$.

    Dessa forma, para cada $p \in P$, obtemos uma função $f_p: \omega\to A$.
    Seja $f:P\times \omega\to A$ dada por $f(p, n)=f_p(n)$.
    Dado $p \in P$ e $n \in \omega$, seja $t$ uma computação parcial baseada em $G$ e $p$ de comprimento $S(n)$.
    Temos que $t\subseteq f_p$, e, portanto, $t\restriction n= f_p\restriction n$.
    Assim, $f_p(n)=t(n)=G(p, t\restriction n)=G(p, f_p\restriction n)$, o que completa a prova da existência.

    Para a unicidade, veremos que $f:P\times \omega\to A$ e $g:P\times \omega\to A$ são tais que, para todo $p \in P$ e $n\in \omega$, $f_p(n)=G(p, f_p\restriction n)$ e $g_p(n)=G(p, g_p\restriction n)$, então $f=g$.

    Fixado $p$, veremos por indução forte que para todo $n\in \omega$, $f_p(n)=g_p(n)$.
    Provada a tese para todo $k<n$, temos que $f_p\restriction n=g_p\restriction n$.
    Assim, $f_p(n)=G(p, f_p\restriction n)=G(p, g_p\restriction n)=g_p(n)$, o que completa a indução.
\end{proof}

\section*{Exercícios}
\begin{exercise}[Teorema da recursão sem parâmetros]
    Prove que para todo conjunto $A$ e para todo $G: A^{<\omega}\to A$, existe uma única função $f:\omega\to A$ tal que, para todo $n\in \omega$, $f(n)=G(f\restriction n)$.
\end{exercise}